\documentclass{beamer}

\title{Week 2 + 3: Policy Gradient Methods}
\author{Phuc Nguyen Dang}
\date{\today}

\begin{document}

\frame{\titlepage}

\begin{frame}{Policy Gradient: Learn a Policy Directly}

\textbf{Objective:} learn a parameterized policy
\[
\pi_\theta(a|s)
\]
by maximizing a performance measure:
\[
J(\theta)
\]

\vspace{0.2cm}

\textbf{Gradient ascent update}
\[
\theta_{t+1}
=
\theta_t + \alpha \nabla_\theta J(\theta_t)
\]

\vspace{0.2cm}

\textbf{Policy Gradient Theorem}
\[
\begin{aligned}
\nabla_\theta J(\theta)
&\propto
\sum_s \mu_\pi(s)
\sum_a
\nabla_\theta \pi_\theta(a|s) Q^\pi(s,a) \\
&=
\sum_s \mu_\pi(s)
\sum_a
\pi_\theta(a|s)
\nabla_\theta \log \pi_\theta(a|s)
Q^\pi(s,a) \\
&=
\mathbb{E}_{s \sim \mu_\pi,\ a \sim \pi_\theta}
\left[
\nabla_\theta \log \pi_\theta(a|s) Q^\pi(s,a)
\right].
\end{aligned}
\]
\end{frame}

\begin{frame}{From Policy Gradient to Algorithms}

From the policy gradient theorem, the gradient can be estimated using sampled trajectories:

\[
\nabla_\theta J(\theta)
\approx
\mathbb{E}_t
\left[
\nabla_\theta \log \pi_\theta(A_t|S_t)
\; \widehat{\Psi}_t
\right]
\]

where \(\widehat{\Psi}_t\) measures how good the sampled action was.

\vspace{0.3cm}

\begin{block}{Two key questions}
\begin{enumerate}
    \item \textbf{How good or bad was the sampled action?}
    \[
    \widehat{\Psi}_t
    =
    G_t,\quad
    G_t - b(S_t),\quad
    A_t,\quad
    \delta_t,\quad
    \widehat{A}^{\mathrm{GAE}}_t
    \]

    \item \textbf{How strong should the policy update be?}
    \[
    \theta_{\text{new}}
    \leftarrow
    \theta_{\text{old}}
    + \text{controlled policy update}
    \]
\end{enumerate}
\end{block}

\vspace{0.2cm}

\begin{center}
REINFORCE \(\rightarrow\) Baseline \(\rightarrow\) Actor-Critic / GAE
\(\rightarrow\) TRPO \(\rightarrow\) PPO
\end{center}

\end{frame}
\begin{frame}{REINFORCE: Policy Gradient Estimator}

\small

\[
\nabla J(\theta)
\propto
\sum_s \mu_\gamma(s)
\sum_a q_\pi(s,a)\nabla_\theta \pi(a|s)
=
\mathbb{E}_{S_t \sim \mu_\gamma}
\left[
\sum_a q_\pi(S_t,a)
\nabla_\theta \pi(a|S_t,\theta)
\right]
\]

\vspace{0.2cm}

\begin{align*}
\nabla J(\theta)
&\propto
\mathbb{E}_{S_t \sim \mu_\gamma}
\left[
\sum_a
\pi(a|S_t,\theta)
q_\pi(S_t,a)
\frac{
\nabla_\theta \pi(a|S_t,\theta)
}{
\pi(a|S_t,\theta)
}
\right]
\tag{1}
\\[0.15cm]
&=
\mathbb{E}_{S_t \sim \mu_\gamma,\ A_t \sim \pi}
\left[
q_\pi(S_t,A_t)
\nabla_\theta \ln \pi(A_t|S_t,\theta)
\right]
\tag{2}
\\[0.15cm]
&=
\mathbb{E}_{S_t \sim \mu_\gamma,\ A_t \sim \pi}
\left[
G_t
\nabla_\theta \ln \pi(A_t|S_t,\theta)
\right],
\tag{3}
\end{align*}

\vspace{0.1cm}

Because \(\mu_\gamma\) weights the state at time \(t\) by \(\gamma^t\),
the Monte-Carlo episode update uses the factor \(\gamma^t\).

\end{frame}

\begin{frame}{REINFORCE: Monte-Carlo Policy-Gradient Control}

\small

\textbf{Goal:} learn a parameterized stochastic policy
\[
\pi(a|s,\theta)
\]

\vspace{0.15cm}

\textbf{Input:}
\[
\pi(a|s,\theta), \qquad \alpha > 0
\]

\vspace{0.15cm}

\textbf{Initialize:}
\[
\theta \in \mathbb{R}^d
\]

\vspace{0.15cm}

\textbf{For each episode:}

\begin{itemize}
    \item Generate one trajectory following the current policy:
    \[
    S_0, A_0, R_1, \ldots, S_{T-1}, A_{T-1}, R_T
    \]

    \item For each time step \(t = 0,1,\ldots,T-1\):
    \[
    G_t
    =
    \sum_{k=t+1}^{T}
    \gamma^{k-t-1} R_k
    \]

    \[
    \theta
    \leftarrow
    \theta
    +
    \alpha \gamma^t G_t
    \nabla_\theta \ln \pi(A_t|S_t,\theta)
    \]
\end{itemize}

\end{frame}

\begin{frame}{REINFORCE with Baseline: Policy Gradient Estimator}

\small

From the Appendix:
\[
\nabla J(\theta)
\propto
\sum_s \mu_\gamma(s)
\sum_a q_\pi(s,a)\nabla_\theta \pi(a|s)
\]

\vspace{0.1cm}

Adding a baseline \(b(s)\) to the action-value function gives:
\[
\nabla J(\theta)
\propto
\sum_s \mu_\gamma(s)
\sum_a
\left(q_\pi(s,a) - b(s)\right)
\nabla_\theta \pi(a|s)
\]

\vspace{0.1cm}

As long as \(b(s)\) does not depend on \(a\), the expected value of
the policy gradient remains unchanged:
\[
\sum_a b(s)\nabla_\theta \pi(a|s)
=
b(s)\nabla_\theta \sum_a \pi(a|s)
=
b(s)\nabla_\theta 1
=
0.
\]

\vspace{0.1cm}

With the baseline, the update rule for REINFORCE becomes:
\[
\theta_{t+1}
\leftarrow
\theta_t
+
\alpha \gamma^t
\left(G_t - b(S_t)\right)
\nabla_\theta \ln \pi(A_t|S_t,\theta)
\]

\end{frame}

\begin{frame}{REINFORCE with Baseline: Monte-Carlo Policy-Gradient Control}

\scriptsize

\begin{tabular}{@{}p{0.24\linewidth}p{0.70\linewidth}@{}}
\textbf{Natural baseline:}
&
\(b(S_t)=\hat{v}(S_t,w),\quad w\in\mathbb{R}^{d}\)
\\[0.08cm]
\textbf{Input:}
&
\(\pi(a|s,\theta),\quad \hat{v}(s,w),\quad
\alpha^\theta>0,\ \alpha^w>0\)
\\[0.08cm]
\textbf{Initialize:}
&
\(\theta\in\mathbb{R}^{d'},\quad w\in\mathbb{R}^{d}\)
\end{tabular}

\vspace{0.12cm}

\textbf{For each episode:}
\begin{enumerate}
    \item Generate one trajectory following the current policy:
    \[
    S_0, A_0, R_1, \ldots, S_{T-1}, A_{T-1}, R_T
    \]

    \item For each time step \(t=0,1,\ldots,T-1\):
    \[
    \begin{aligned}
    G_t
    &=
    \sum_{k=t+1}^{T}
    \gamma^{k-t-1} R_k,
    &
    \delta
    &=
    G_t - \hat{v}(S_t,w)
    \end{aligned}
    \]

    \[
    \begin{aligned}
    w
    &\leftarrow
    w
    +
    \alpha^w \delta
    \nabla_w \hat{v}(S_t,w),
    \\
    \theta
    &\leftarrow
    \theta
    +
    \alpha^\theta \gamma^t \delta
    \nabla_\theta \ln \pi(A_t|S_t,\theta)
    \end{aligned}
    \]
\end{enumerate}

\end{frame}

\begin{frame}{From REINFORCE with Baseline to Actor-Critic}

\small

\textbf{Remaining issue: still Monte-Carlo.}

\vspace{0.1cm}

Even with a baseline, REINFORCE still needs the return:
\[
G_t
=
\sum_{k=t+1}^{T}
\gamma^{k-t-1}R_k
\]

\vspace{0.1cm}

Thus, the return is known only after the episode finishes.

\vspace{0.15cm}

\begin{block}{Three practical issues}
\begin{itemize}
    \item The episode must finish before \(G_t\) is available.
    \item \(G_t\) can still be noisy.
    \item Online learning and continuing tasks become harder.
\end{itemize}
\end{block}

\vspace{0.1cm}

\[
\text{Actor-Critic replaces } G_t
\text{ with a bootstrapped TD target.}
\]

\end{frame}

\begin{frame}{Actor-Critic: From Monte-Carlo to TD}

\small

TD error stands for Temporal-Difference error.
It measures how wrong the value prediction was after observing one more step.

\[
\delta_t
=
R_{t+1}
+
\gamma \hat{v}(S_{t+1},w_t)
-
\hat{v}(S_t,w_t)
\]

\vspace{0.1cm}

Instead of using the Monte-Carlo return \(G_t\) to update the policy,
Actor-Critic uses the one-step TD target:
\[
R_{t+1}
+
\gamma \hat{v}(S_{t+1},w)
\]

\vspace{0.1cm}

\begin{itemize}
    \item If \(\hat{v}\) is accurate, the TD target estimates \(q_\pi(S_t,A_t)\).
    \item In practice, the target is biased but usually has lower variance than \(G_t\).
\end{itemize}

\end{frame}

\begin{frame}{Actor-Critic: Actor, Critic, and Signal}

\small

\begin{itemize}
    \item \textbf{Actor:} the policy \(\pi(a|s,\theta)\)
    \item \textbf{Critic:} the state-value estimate \(\hat{v}(s,w)\)
    \item \textbf{Signal:} the TD error
\end{itemize}

\[
\delta
\leftarrow
R
+
\gamma \hat{v}(S',w)
-
\hat{v}(S,w)
\]

\vspace{0.1cm}

The TD error plays the role of the advantage estimate used to update both
the critic and the actor.

\[
w
\leftarrow
w
+
\alpha^w \delta
\nabla_w \hat{v}(S,w)
\]

\[
\theta
\leftarrow
\theta
+
\alpha^\theta \delta I
\nabla_\theta \ln \pi(A|S,\theta)
\]

where \(I=\gamma^t\) in episodic policy-gradient updates.

\end{frame}

\begin{frame}{Forward View: \(n\)-Step and \(\lambda\)-Return}

\small

The one-step return can be replaced by an \(n\)-step target:
\[
G_{t:t+n}
=
R_{t+1}
+
\gamma R_{t+2}
+
\cdots
+
\gamma^{n-1}R_{t+n}
+
\gamma^n \hat{v}(S_{t+n},w)
\]

\vspace{0.1cm}

The \(\lambda\)-return is a weighted average of these \(n\)-step targets:
\[
G_t^\lambda
=
(1-\lambda)
\sum_{n=1}^{T-t-1}
\lambda^{n-1}G_{t:t+n}
+
\lambda^{T-t-1}G_{t:T}
\]

\vspace{0.1cm}

\[
\lambda=0:\ \text{one-step TD target},
\qquad
\lambda=1:\ \text{Monte-Carlo return}
\]

Trade off: bias and variance.
\end{frame}

\begin{frame}{Forward View: Actor-Critic Update}

\scriptsize

For the critic, the return error is a prediction error:
\[
e_t^\lambda
=
G_t^\lambda
-
\hat{v}(S_t,w_t)
\]

\vspace{0.05cm}

The critic is updated by semi-gradient descent:
\[
w_{t+1}
=
w_t
+
\alpha^w e_t^\lambda
\nabla_w \hat{v}(S_t,w_t)
\]

\[
-\nabla_w
\frac{1}{2}
\left(G_t^\lambda-\hat{v}(S_t,w_t)\right)^2
=
e_t^\lambda \nabla_w \hat{v}(S_t,w_t)
\]

\vspace{0.05cm}

For the actor, the same error acts like an advantage estimate:
\[
\widehat{A}_t^\lambda
\doteq
e_t^\lambda
=
G_t^\lambda
-
\hat{v}(S_t,w_t)
\]

The actor update is:
\[
\theta_{t+1}
=
\theta_t
+
\alpha^\theta \gamma^t
\widehat{A}_t^\lambda
\nabla_\theta \ln \pi(A_t|S_t,\theta_t)
\]

\end{frame}

\begin{frame}{Backward View: Eligibility Traces}

\small

Computing \(G_t^\lambda\) directly is the forward view because it needs
future rewards.

\vspace{0.1cm}

The backward view implements the same idea online with eligibility traces:
\[
G_t^\lambda - \hat{v}(S_t,w)
\approx
\sum_{k=t}^{T-1}
(\gamma\lambda)^{k-t}
\delta_k
\]

\vspace{0.1cm}

Critic trace:
\[
z_t^w
=
\gamma \lambda^w z_{t-1}^w
+
\nabla_w \hat{v}(S_t,w_t)
\]

Actor trace:
\[
z_t^\theta
=
\gamma \lambda^\theta z_{t-1}^\theta
+
I_t \nabla_\theta \ln \pi(A_t|S_t,\theta_t)
\]

\end{frame}

\begin{frame}{Actor-Critic with Eligibility Traces}

\scriptsize

\textbf{Initialize each episode:}
\[
S,\qquad
z^\theta \leftarrow \mathbf{0},
\qquad
z^w \leftarrow \mathbf{0},
\qquad
I \leftarrow 1
\]

\vspace{0.1cm}

\textbf{For each step while \(S\) is not terminal:}
\begin{itemize}
    \item \(A \sim \pi(\cdot|S,\theta)\); take \(A\), observe \(S'\), \(R\)
    \item
    \[
    \delta
    \leftarrow
    R
    +
    \gamma \hat{v}(S',w)
    -
    \hat{v}(S,w)
    \]
    \item
    \[
    z^w
    \leftarrow
    \gamma\lambda^w z^w
    +
    \nabla_w \hat{v}(S,w)
    \]
    \[
    z^\theta
    \leftarrow
    \gamma\lambda^\theta z^\theta
    +
    I\nabla_\theta\ln\pi(A|S,\theta)
    \]
\end{itemize}

\[
w
\leftarrow
w
+
\alpha^w \delta z^w,
\qquad
\theta
\leftarrow
\theta
+
\alpha^\theta \delta z^\theta
\]

\[
I \leftarrow \gamma I,
\qquad
S \leftarrow S'
\]

\end{frame}


\begin{frame}{Second Question: How Strong Should the Update Be?}

\small

So far, the algorithms mainly answer the first question:
\[
\text{How good or bad was the sampled action?}
\]

\[
G_t,\quad
G_t-b(S_t),\quad
\delta_t,\quad
\widehat{A}_t^{\lambda}
\]

\vspace{0.15cm}

Given an advantage estimate \(\widehat{A}_t\), the next question is:
\[
\text{How much should the policy } \pi_\theta \text{ change?}
\]

\begin{block}{Problem}
Training data is sampled from the old policy
\(\pi_{\theta_{\mathrm{old}}}\),
but used to update the new policy \(\pi_\theta\).
If the policy changes too much, the old data may no longer be useful.
\end{block}

\[
\text{A controlled policy update is needed.}
\]

\end{frame}

\begin{frame}{Performance Difference Between Two Policies}

\small

Let \(\pi\) be the old policy and \(\tilde{\pi}\) be the new policy.
Using the advantage under the old policy:
\[
J(\tilde{\pi})
=
J(\pi)
+
\mathbb{E}_{\tilde{\pi}}
\left[
\sum_{t=0}^{\infty}
\gamma^t A^\pi(s_t,a_t)
\right]
\]

\[
=
J(\pi)
+
\sum_s
\mu_{\tilde{\pi}}(s)
\sum_a
\tilde{\pi}(a|s)
A^\pi(s,a)
\]

\vspace{0.1cm}

This says:
\[
\text{new performance}
=
\text{old performance}
\]
\[
+
\text{expected old-policy advantage under the new policy.}
\]

\begin{block}{Difficulty}
\(\mu_{\tilde{\pi}}(s)\) also depends on the new policy,
so directly optimizing this expression is hard.
\end{block}

\end{frame}

\begin{frame}{Surrogate Objective: Reuse Old-Policy Data}

\small

To make optimization practical, approximate the new state distribution
by the old one:
\[
\mu_{\tilde{\pi}}(s) \approx \mu_{\pi}(s)
\]

\[
L^{\mathrm{CPI}}(\tilde{\pi})
=
J(\pi)
+
\sum_s
\mu_{\pi}(s)
\sum_a
\tilde{\pi}(a|s)
A^\pi(s,a)
\]

\vspace{0.1cm}

Rewrite it as an expectation over old-policy samples:
\[
\begin{aligned}
L^{\mathrm{CPI}}(\tilde{\pi}) - J(\pi)
&=
\sum_s \mu_{\pi}(s)
\sum_a
\pi(a|s)
\frac{\tilde{\pi}(a|s)}{\pi(a|s)}
A^\pi(s,a)
\\
&=
\mathbb{E}_{s,a\sim\pi}
\left[
\frac{\tilde{\pi}(a|s)}{\pi(a|s)}
A^\pi(s,a)
\right].
\end{aligned}
\]

\end{frame}

\begin{frame}{TRPO: Trust-Region Policy Update}

\small

With parameters, the surrogate objective becomes:
\[
\max_{\theta}
\mathbb{E}_{s,a\sim\pi_{\theta_{\mathrm{old}}}}
\left[
\frac{
\pi_{\theta}(a|s)
}{
\pi_{\theta_{\mathrm{old}}}(a|s)
}
A^{\pi_{\theta_{\mathrm{old}}}}(s,a)
\right]
\]

\vspace{0.1cm}

But the approximation is valid only when the new policy stays close
to the old policy.

\[
D_{\mathrm{KL}}
\left(
\pi_{\theta_{\mathrm{old}}}
\;\|\;
\pi_{\theta}
\right)
\leq
\delta
\]

\begin{block}{TRPO idea}
Maximize the surrogate improvement, but constrain the KL divergence
between old and new policies.
\end{block}

\[
\delta
=
\text{the maximum reasonable policy update size.}
\]

\end{frame}

\begin{frame}{From TRPO to PPO}

\small

TRPO gives a principled answer:
\[
\text{update enough to improve the surrogate,}
\qquad
\text{but not beyond a KL trust region.}
\]

\vspace{0.15cm}

However, this is a constrained optimization problem:
\[
\max_\theta L(\theta)
\quad
\text{subject to}
\quad
D_{\mathrm{KL}}
\left(
\pi_{\theta_{\mathrm{old}}}
\|\pi_\theta
\right)
\leq \delta
\]

\vspace{0.15cm}

Solving it usually requires second-order approximations and conjugate
gradient steps.

\begin{block}{PPO idea}
Replace the hard KL constraint with a simpler clipped objective.
\end{block}

\[
\text{The update size is controlled by clipping the probability ratio.}
\]

\end{frame}

\begin{frame}{PPO: Clipped Surrogate Objective}

\small

Define the probability ratio:
\[
r_t(\theta)
=
\frac{
\pi_{\theta}(A_t|S_t)
}{
\pi_{\theta_{\mathrm{old}}}(A_t|S_t)
}
\]

\vspace{0.1cm}

PPO maximizes:
\[
L^{\mathrm{CLIP}}(\theta)
=
\mathbb{E}_t
\left[
\min
\left(
r_t(\theta)\widehat{A}_t,\;
\operatorname{clip}
\left(
r_t(\theta), 1-\epsilon, 1+\epsilon
\right)
\widehat{A}_t
\right)
\right]
\]

\vspace{0.1cm}


\end{frame}

\begin{frame}{PPO Clipping: Single Timestep}

\centering

\IfFileExists{ppo_clip_term.png}{
    \includegraphics[
        width=0.95\linewidth,
        height=0.72\textheight,
        keepaspectratio
    ]{ppo_clip_term.png}
}{
    \fbox{
        \parbox[c][0.72\textheight][c]{0.88\linewidth}{
            \centering
            \texttt{ppo\_clip\_term.png}
        }
    }
}

\vspace{0.05cm}
{\scriptsize Clip term}

\end{frame}

\begin{frame}{PPO Clipping: Objectives}

\centering

\IfFileExists{ppo_objectives.png}{
    \includegraphics[
        width=0.95\linewidth,
        height=0.72\textheight,
        keepaspectratio
    ]{ppo_objectives.png}
}{
    \fbox{
        \parbox[c][0.72\textheight][c]{0.88\linewidth}{
            \centering
            \texttt{ppo\_objectives.png}
        }
    }
}

\vspace{0.05cm}
{\scriptsize Objectives}

\end{frame}

\end{document}
